% Preambuła zgodna ze skillem latex-engineer (.claude/skills/latex-engineer):
% pakiety pogrupowane funkcyjnie, jeden plik, bez pakietów w main.tex.

% --- Język (PL) ---
\usepackage[T1]{fontenc}
\usepackage[utf8]{inputenc}
\usepackage[polish]{babel}
\usepackage{indentfirst}

% --- Typografia i układ strony ---
\usepackage{lmodern}
\usepackage{microtype}
\usepackage[all]{nowidow}
\usepackage[a4paper,margin=2.5cm]{geometry}
\usepackage{needspace}
\usepackage{enumitem}

% --- Grafika i tabele ---
\usepackage{graphicx}
\usepackage{booktabs}
\usepackage{tabularx}
\usepackage{caption}
\usepackage{xcolor}

% --- Kod i listingi ---
\usepackage{listings}

% --- Nagłówek i stopka ---
\usepackage{fancyhdr}
\usepackage{lastpage}
\pagestyle{fancy}
\fancyhf{}
\lhead{\leftmark}
\rhead{RAG: Asystent Sylabusów}
\cfoot{Strona \thepage{} z \pageref{LastPage}}
\renewcommand{\headrulewidth}{0.4pt}

% --- Bibliografia i linki ---
\usepackage[backend=bibtex,style=numeric,sorting=none]{biblatex}
\addbibresource{bib/zrodla.bib}
\usepackage[colorlinks=true,linkcolor=blue,urlcolor=blue,citecolor=blue]{hyperref}

% --- Makra projektu ---
% Listingi z polskimi znakami (listings ignoruje inputenc)
\lstdefinestyle{prompt}{
  basicstyle=\ttfamily\small,
  breaklines=true,
  breakatwhitespace=true,
  columns=fullflexible,
  frame=single,
  rulecolor=\color{black!30},
  backgroundcolor=\color{black!3},
  xleftmargin=0.5em,
  xrightmargin=0.5em,
  extendedchars=true,
  inputencoding=utf8,
  literate=
    {ą}{{\k{a}}}1 {ć}{{\'c}}1 {ę}{{\k{e}}}1 {ł}{{\l{}}}1
    {ń}{{\'n}}1 {ó}{{\'o}}1 {ś}{{\'s}}1 {ź}{{\'z}}1 {ż}{{\.z}}1
    {Ą}{{\k{A}}}1 {Ć}{{\'C}}1 {Ę}{{\k{E}}}1 {Ł}{{\L{}}}1
    {Ń}{{\'N}}1 {Ó}{{\'O}}1 {Ś}{{\'S}}1 {Ź}{{\'Z}}1 {Ż}{{\.Z}}1
}
\lstset{style=prompt}

% Wpis dziennika promptow: #1 nr, #2 data, #3 model, #4 effort, #5 cel.
% Tresc promptu w srodowisku lstlisting bezposrednio po makrze,
% a po nim \promptresult{...} z krotkim opisem wyniku.
\newcommand{\promptentry}[5]{%
  \needspace{8\baselineskip}%
  \subsection*{Prompt #1}%
  \begin{tabularx}{\linewidth}{@{}lX@{}}
    \toprule
    Data & #2 \\
    Model & #3 \\
    Effort & #4 \\
    Cel & #5 \\
    \bottomrule
  \end{tabularx}\par\smallskip
}
\newcommand{\promptresult}[1]{\par\smallskip\noindent\textbf{Wynik:} #1\par\bigskip}

% Oznaczenie tresci wygenerowanych przez AI
\newcommand{\aigen}[1]{#1\textsuperscript{\textsc{ai}}}

